\section{Summary and Future Work}
\label{sec:summary-future-work}

The BiteBound project demonstrates a practical approach to addressing the primary objectives established in~\Cref{sec:introduction}, specifically managing low-latency gameplay execution alongside secure, asynchronous network communication on resource-constrained hardware.
By partitioning the timing-sensitive physics engine and rendering pipeline onto the Game Core and the communication onto the Network Core (Core 0), the system maintains a stable frame rate.
This separation of concerns prevents network latency from interrupting physical gameplay while maintaining bidirectional synchronization with the Android and Node-RED companion interfaces, validating the multi-platform user experience proposed in the mission statement.

Future developments will focus on expanding the platform's physical interactivity, sensory feedback, and academic dissemination.
To enrich the gameplay mechanics, upcoming iterations can incorporate onboard physical button inputs to trigger special, context-dependent actions during active play.
Integrating a vibration motor or haptic sensor would further ground the physical interaction by providing physical feedback when the virtual ball collides with maze boundaries.
Auditory feedback can also be expanded by implementing driver routines for the onboard buzzer, generating custom startup sequences, loading music, and distinctive alerts for coin-collection events.
Finally, this technical report will be formatted into a scientific paper suitable for public submission to relevant mobile and ubiquitous computing venues, allowing for broader academic evaluation of the architecture.